\documentclass[UTF8,a4paper,zihao=-4,fontset=fandol]{ctexart}
\usepackage{wefjnu-exam}

\begin{document}
\wefjnukeytitle{《高等数学 C（上）》期末考试参考答案与解析（A 卷）}

% 只填原卷中可以核实的信息。只有一两项时可用 \wefjnuinfo 单独排一行；
% 没有可核实的信息时，省略这些信息行。
\wefjnumeta{适用专业　全校各相关专业}{适用年级　2025 级}
  {考试时间　2026 年 1 月 14 日 09:00}{任课教师　朱金才等}
\wefjnuinfo{闭卷 \quad 120 分钟 \quad 满分 100 分}

% 将「实际模型名」替换为该答案实际使用的模型。
\wefjnuaisource{实际模型名}

% 按原卷题型与题号编排，给出结论和关键推导。
\wefjnusection{一、选择题}
\begin{enumerate}
  \item \wefjnuanswer{A（可去间断点）。}
  \wefjnuwork
  原函数在 $x=0$ 处无定义。因为 $\tan x/x\to1$，
  所以 $\displaystyle\lim_{x\to0}x/\tan x=1$，补定义 $f(0)=1$ 后连续。
\end{enumerate}
\end{document}
